\chapter{Goniometrie: de eenheidscirkel}\label{ch:g11:trigo}

De goniometrie begint in rechthoekige driehoeken, maar haar echte thuis is de
\emph{eenheidscirkel}, waar \omterm{def:g11:trigo:cossin}{cosinus} en \omterm{def:g11:trigo:cossin}{sinus} \omterm{def:g11:func:function}{functies} van een willekeurig
\omterm{def:g10:numbers:sets}{reëel getal} worden. Dit hoofdstuk installeert de \omterm{def:g11:trigo:radian}{radiaal}, het oprollen van de
reële lijn om de cirkel, de exacte waarden die je uit het hoofd moet kennen,
en de symmetrieën die alle klassieke identiteiten voortbrengen. De studie van
$\cos$ en $\sin$ als \omterm{def:g11:func:function}{functies} --- verloop, \omterm{def:g11:deriv:derivative}{afgeleiden} --- gebeurt in
\cref{ch:g12:trigo}.

\section{Radialen en het oprollen}

\begin{definition}[Radiaal]\label{def:g11:trigo:radian}
Zij $\mathcal C$ de cirkel met straal $1$ en middelpunt de oorsprong $O$, en
$I(1, 0)$. De maat in \emph{radialen}\index{radiaal} van een middelpuntshoek
van $\mathcal C$ is de lengte van de boog die hij op $\mathcal C$ afsnijdt.
Omdat de volledige cirkel omtrek $2\pi$ heeft, geldt
\[
360^\circ = 2\pi \text{ rad},
\qquad
180^\circ = \pi \text{ rad},
\qquad
1 \text{ rad} = \frac{180^\circ}{\pi} \approx 57.3^\circ .
\]
\end{definition}

\begin{method}[Omrekenen]\label{met:g11:trigo:convert}
Graden en \omterm{def:g11:trigo:radian}{radialen} zijn evenredig: vermenigvuldig met $\frac{\pi}{180}$ om
van graden naar \omterm{def:g11:trigo:radian}{radialen} te gaan, en met $\frac{180}{\pi}$ in de andere
richting. Zo is $60^\circ = 60 \times \frac{\pi}{180} = \frac{\pi}{3}$ en
$\frac{3\pi}{4} = \frac{3 \times 180^\circ}{4} = 135^\circ$.
\end{method}

\begin{definition}[De lijn om de cirkel oprollen]\label{def:g11:trigo:winding}
Koppel aan elk \omterm{def:g10:numbers:sets}{reëel getal} $t$ het punt $M(t)$ van $\mathcal C$ dat je
bereikt door vanaf $I$ een afstand $\abs t$ langs de cirkel af te leggen,
tegen de wijzers van de klok in wanneer $t \geq 0$ en met de wijzers mee
wanneer $t < 0$. Omdat de omtrek $2\pi$ is, bereiken de getallen $t$ en
$t + 2k\pi$ ($k \in \Z$) hetzelfde punt: $M(t + 2k\pi) = M(t)$.
\end{definition}

\begin{example}
$M(0) = I$; $M\!\left(\frac{\pi}{2}\right) = (0, 1)$, de top van de cirkel;
$M(\pi) = (-1, 0)$; $M\!\left(-\frac{\pi}{2}\right) = (0, -1)$; en
$M\!\left(\frac{9\pi}{4}\right) = M\!\left(\frac{\pi}{4} + 2\pi\right) =
M\!\left(\frac{\pi}{4}\right)$.
\end{example}

\section{Cosinus en sinus}

\begin{definition}[Cosinus en sinus]\label{def:g11:trigo:cossin}
Voor elke reële $t$ zijn de \emph{cosinus}\index{cosinus} en de
\emph{sinus}\index{sinus} van $t$ de \omterm{def:g10:coordgeom:system}{coördinaten} van het punt $M(t)$ van de
eenheidscirkel:
\[
M(t) = (\cos t,\ \sin t).
\]
\end{definition}

\begin{omfigure}
\begin{tikzpicture}[scale=2.0]
  \draw[->] (-1.3,0) -- (1.4,0) node[right] {$x$};
  \draw[->] (0,-1.3) -- (0,1.4) node[above] {$y$};
  \draw[gray] (0,0) circle (1);
  \coordinate (O) at (0,0);
  \coordinate (M) at (55:1);
  \draw[thick,omDef] (O) -- (M);
  \draw[dashed] (M) -- (M |- O) node[below, font=\small] {$\cos t$};
  \draw[dashed] (M) -- (M -| O) node[left, font=\small] {$\sin t$};
  \draw[->,omThm,thick] (0.35,0) arc (0:55:0.35);
  \node[omThm, font=\small] at (27:0.55) {$t$};
  \fill (M) circle (0.7pt) node[above right] {$M(t)$};
  \fill (1,0) circle (0.7pt) node[below right] {$I$};
\end{tikzpicture}

{\small De eenheidscirkel: het reële getal $t$, vanaf $I$ tegen de wijzers
van de klok in opgerold, komt uit in het punt $M(t)$, waarvan de \omterm{def:g10:coordgeom:system}{coördinaten}
$\cos t$ en $\sin t$ \emph{definiëren}.}
\end{omfigure}

\begin{proposition}[Eerste eigenschappen]\label{prop:g11:trigo:first}
Voor alle $t \in \R$ en $k \in \Z$ geldt
\[
-1 \leq \cos t \leq 1, \qquad
-1 \leq \sin t \leq 1, \qquad
\cos^2 t + \sin^2 t = 1,
\]
\[
\cos(t + 2k\pi) = \cos t, \qquad \sin(t + 2k\pi) = \sin t .
\]
\end{proposition}

\begin{proof}
$M(t)$ ligt op de eenheidscirkel, dus liggen zijn \omterm{def:g10:coordgeom:system}{coördinaten} tussen $-1$ en
$1$, en voldoen ze aan de \omterm{def:g10:algebra:equation}{vergelijking} $x^2 + y^2 = 1$ van de cirkel, wat de
identiteit $\cos^2 t + \sin^2 t = 1$ is. De periodiciteit herhaalt
$M(t + 2k\pi) = M(t)$ (\cref{def:g11:trigo:winding}).
\end{proof}

\begin{example}\label{ex:g11:trigo:identity}
Is $\sin t = \frac35$ en $t \in \intcc{\frac\pi2}{\pi}$ (tweede kwadrant),
dan is $\cos^2 t = 1 - \frac{9}{25} = \frac{16}{25}$, dus
$\cos t = \pm\frac45$; in het tweede kwadrant is de \omterm{def:g10:coordgeom:system}{abscis} negatief, en dus
is $\cos t = -\frac45$.
\end{example}

De waarden die je uit het hoofd moet kennen:
\begin{center}
\begin{tabular}{c|ccccc}
$t$ & $0$ & $\dfrac{\pi}{6}$ & $\dfrac{\pi}{4}$ & $\dfrac{\pi}{3}$ & $\dfrac{\pi}{2}$\\[6pt]
\hline\rule{0pt}{14pt}%
$\cos t$ & $1$ & $\dfrac{\sqrt3}{2}$ & $\dfrac{\sqrt2}{2}$ & $\dfrac12$ & $0$\\[6pt]
$\sin t$ & $0$ & $\dfrac12$ & $\dfrac{\sqrt2}{2}$ & $\dfrac{\sqrt3}{2}$ & $1$
\end{tabular}
\end{center}

\begin{remark}
Een ezelsbruggetje: de sinussen lezen
$\frac{\sqrt0}{2}, \frac{\sqrt1}{2}, \frac{\sqrt2}{2}, \frac{\sqrt3}{2},
\frac{\sqrt4}{2}$, en de cosinussen zijn dezelfde lijst omgekeerd.
\end{remark}

\section{Verwante hoeken}

\begin{proposition}[Verwante hoeken]\label{prop:g11:trigo:associated}
\index{verwante hoeken}
Voor alle $t \in \R$ geldt
\[
\begin{aligned}
\cos(-t) &= \cos t, & \sin(-t) &= -\sin t,\\
\cos(\pi - t) &= -\cos t, & \sin(\pi - t) &= \sin t,\\
\cos(\pi + t) &= -\cos t, & \sin(\pi + t) &= -\sin t,\\
\cos\left(\tfrac{\pi}{2} - t\right) &= \sin t, &
\sin\left(\tfrac{\pi}{2} - t\right) &= \cos t .
\end{aligned}
\]
\end{proposition}

\begin{proof}
Elke regel drukt een symmetrie van de cirkel uit.

$M(-t)$ is het spiegelbeeld van $M(t)$ om de $x$-as: de \omterm{def:g10:coordgeom:system}{abscis} blijft, de
\omterm{def:g10:coordgeom:system}{ordinaat} wisselt van teken.

$M(\pi - t)$ is het spiegelbeeld van $M(t)$ om de $y$-as: vanaf $\pi$ over
$t$ terugwandelen komt horizontaal tegenover vanaf $0$ over $t$
vooruitwandelen uit. De \omterm{def:g10:coordgeom:system}{abscis} wisselt van teken, de \omterm{def:g10:coordgeom:system}{ordinaat} blijft.

$M(\pi + t)$ ligt diametraal tegenover $M(t)$ (een halve draai): beide
\omterm{def:g10:coordgeom:system}{coördinaten} wisselen van teken.

$M\!\left(\frac\pi2 - t\right)$ is het spiegelbeeld van $M(t)$ om de
diagonale rechte $y = x$, die de twee \omterm{def:g10:coordgeom:system}{coördinaten} verwisselt.
\end{proof}

\begin{omfigure}
\begin{tikzpicture}[scale=2.0]
  \draw[->] (-1.35,0) -- (1.45,0) node[right] {$x$};
  \draw[->] (0,-1.35) -- (0,1.45) node[above] {$y$};
  \draw[gray] (0,0) circle (1);
  \fill (40:1) circle (0.7pt) node[above right, font=\small] {$M(t)$};
  \fill (-40:1) circle (0.7pt) node[below right, font=\small] {$M(-t)$};
  \fill (140:1) circle (0.7pt) node[above left, font=\small]
    {$M(\pi - t)$};
  \fill (220:1) circle (0.7pt) node[below left, font=\small]
    {$M(\pi + t)$};
  \draw[black, dashed, thin] (40:1) -- (-40:1);
  \draw[black, dashed, thin] (40:1) -- (140:1);
  \draw[black, dashed, thin] (140:1) -- (220:1);
  \draw[omDef, thin] (0,0) -- (40:1);
  \draw[omDef, thin] (0,0) -- (140:1);
  \draw[omDef, thin] (0,0) -- (220:1);
  \draw[omDef, thin] (0,0) -- (-40:1);
\end{tikzpicture}

{\small De vier verwante punten: $M(-t)$ spiegelt $M(t)$ om de $x$-as,
$M(\pi - t)$ om de $y$-as, en $M(\pi + t)$ ligt diametraal tegenover. Elke
identiteit van \cref{prop:g11:trigo:associated} lees je van deze tekening
af.}
\end{omfigure}

\begin{example}
$\cos\frac{2\pi}{3} = \cos\left(\pi - \frac\pi3\right) = -\cos\frac\pi3
= -\frac12$, en $\sin\frac{7\pi}{6} = \sin\left(\pi + \frac\pi6\right)
= -\sin\frac\pi6 = -\frac12$. De tabel met vijf waarden, uitgebreid met de
symmetrieën, bestrijkt de hele cirkel.
\end{example}

\section{Goniometrische vergelijkingen oplossen}

\begin{theorem}[De vergelijkingen $\cos t = \cos a$ en $\sin t = \sin a$]
\label{thm:g11:trigo:equations}
Voor \omterm{def:g10:numbers:sets}{reële getallen} $t$ en $a$ geldt
\[
\cos t = \cos a \iff t = a + 2k\pi \ \text{of}\ t = -a + 2k\pi
\quad (k \in \Z),
\]
\[
\sin t = \sin a \iff t = a + 2k\pi \ \text{of}\ t = \pi - a + 2k\pi
\quad (k \in \Z).
\]
\end{theorem}

\begin{proof}
Twee punten van de eenheidscirkel hebben precies dan dezelfde \omterm{def:g10:coordgeom:system}{abscis} wanneer
ze gelijk zijn of elkaars spiegelbeeld om de $x$-as; volgens
\cref{prop:g11:trigo:associated} is het spiegelbeeld van $M(a)$ het punt
$M(-a)$. Dus betekent $\cos t = \cos a$ dat $M(t) = M(a)$ of
$M(t) = M(-a)$, dus $t = \pm a$ op een veelvoud van $2\pi$ na. Net zo
betekenen gelijke ordinaten gelijke punten of spiegelbeelden om de $y$-as,
en het spiegelbeeld van $M(a)$ is $M(\pi - a)$.
\end{proof}

\begin{method}[$\cos t = c$ oplossen op een interval]\label{met:g11:trigo:solve}
\begin{enumerate}
  \item Zoek \emph{één} hoek $a$ met $\cos a = c$, uit de tabel van bekende
        waarden.
  \item Schrijf de algemene oplossingen $t = \pm a + 2k\pi$ op
        (\cref{thm:g11:trigo:equations}).
  \item Kies de \omterm{def:g10:numbers:sets}{gehele getallen} $k$ die in het gevraagde \omterm{def:g10:numbers:interval}{interval} belanden,
        en noem de oplossingen.
\end{enumerate}
Ga voor $\sin t = c$ op dezelfde manier te werk, met $t = a + 2k\pi$ of
$t = \pi - a + 2k\pi$.
\end{method}

\begin{example}\label{ex:g11:trigo:solve}
Los $\cos t = \frac12$ op op $\intoc{-\pi}{\pi}$. Eén oplossing van
$\cos a = \frac12$ is $a = \frac{\pi}{3}$. De algemene oplossingen zijn
$t = \frac\pi3 + 2k\pi$ en $t = -\frac\pi3 + 2k\pi$. Binnen
$\intoc{-\pi}{\pi}$ draagt alleen $k = 0$ bij:
$t \in \left\{-\frac\pi3,\ \frac\pi3\right\}$.

Los $\sin t = -\frac{\sqrt2}{2}$ op op $\intco{0}{2\pi}$. Eén hoek is
$a = -\frac\pi4$; de oplossingen zijn $t = -\frac\pi4 + 2k\pi$ en
$t = \pi + \frac\pi4 + 2k\pi = \frac{5\pi}{4} + 2k\pi$. In
$\intco{0}{2\pi}$: $t = \frac{7\pi}{4}$ (uit $k = 1$) en
$t = \frac{5\pi}{4}$.
\end{example}

\begin{omfigure}
\begin{tikzpicture}[scale=2.0]
  \draw[->] (-1.35,0) -- (1.45,0) node[right] {$x$};
  \draw[->] (0,-1.2) -- (0,1.3) node[above] {$y$};
  \draw[gray] (0,0) circle (1);
  \draw[omProp, thick] (0.5,-1.05) -- (0.5,1.15)
    node[above, font=\small] {$x = \tfrac12$};
  \fill (60:1) circle (0.7pt)
    node[above right, font=\small] {$M\!\left(\tfrac\pi3\right)$};
  \fill (-60:1) circle (0.7pt)
    node[below right, font=\small] {$M\!\left(-\tfrac\pi3\right)$};
  \draw[omDef, thin] (0,0) -- (60:1);
  \draw[omDef, thin] (0,0) -- (-60:1);
\end{tikzpicture}

{\small $\cos t = \frac12$ oplossen: de verticale rechte $x = \frac12$
(oranje) snijdt de eenheidscirkel in twee punten die symmetrisch liggen om
de $x$-as --- vandaar de twee families oplossingen
$t = \pm\frac\pi3 + 2k\pi$.}
\end{omfigure}

\section{Oefeningen}

\begin{exercise}[$\star$]\label{exo:g11:trigo:1}
Zet om naar \omterm{def:g11:trigo:radian}{radialen}: $30^\circ$, $45^\circ$, $120^\circ$, $270^\circ$. Zet
om naar graden: $\frac{\pi}{5}$, $\frac{5\pi}{6}$, $\frac{7\pi}{4}$,
$\frac{2\pi}{9}$.
\end{exercise}

\begin{exercise}[$\star$]\label{exo:g11:trigo:2}
Zet op de eenheidscirkel de punten $M(t)$ uit voor
\[
t = \frac{2\pi}{3},\quad
t = -\frac{\pi}{4},\quad
t = \frac{17\pi}{6},\quad
t = -\frac{7\pi}{2},
\]
nadat je elke waarde modulo $2\pi$ tot een waarde in $\intoc{-\pi}{\pi}$
hebt herleid.
\end{exercise}

\begin{exercise}[$\star$]\label{exo:g11:trigo:3}
Geef met de tabel en de \omterm{prop:g11:trigo:associated}{verwante hoeken} de exacte waarden van
\[
\cos\frac{3\pi}{4}, \quad
\sin\frac{5\pi}{6}, \quad
\cos\left(-\frac{\pi}{3}\right), \quad
\sin\frac{4\pi}{3}, \quad
\cos\frac{11\pi}{6} .
\]
\end{exercise}

\begin{exercise}[$\star$]\label{exo:g11:trigo:4}
Gegeven $\cos t = \frac{5}{13}$ met $t \in \intoo{-\frac\pi2}{0}$, bereken
$\sin t$ exact.
\end{exercise}

\begin{exercise}[$\star\star$]\label{exo:g11:trigo:5}
Vereenvoudig, voor elke reële $x$:
\[
A = \cos(\pi + x) + \cos(-x) + \sin\left(\frac\pi2 - x\right) - \cos x,
\qquad
B = \sin(\pi - x) + \sin(\pi + x) + \sin(-x) .
\]
\end{exercise}

\begin{exercise}[$\star\star$]\label{exo:g11:trigo:6}
Los op op $\intoc{-\pi}{\pi}$:
\[
\cos t = \frac{\sqrt3}{2}, \qquad
\sin t = \frac12, \qquad
\cos t = -1 .
\]
\end{exercise}

\begin{exercise}[$\star\star$]\label{exo:g11:trigo:7}
Los op op $\intco{0}{2\pi}$:
\[
\sin t = -\frac{\sqrt3}{2}, \qquad
\cos t = 0, \qquad
2\sin t + 1 = 0 .
\]
\end{exercise}

\begin{exercise}[$\star\star$]\label{exo:g11:trigo:8}
Los op $\R$ de \omterm{def:g10:algebra:equation}{vergelijking} $\cos 2t = \cos t$ op. (Gebruik
\cref{thm:g11:trigo:equations} met $a = t$, en bespreek de twee families
oplossingen.)
\end{exercise}

\begin{exercise}[$\star\star$]\label{exo:g11:trigo:9}
Toon aan dat voor elke reële $x$ geldt:
\[
\bigl(\cos x + \sin x\bigr)^2 + \bigl(\cos x - \sin x\bigr)^2 = 2 .
\]
\end{exercise}

\begin{exercise}[$\star\star$]\label{exo:g11:trigo:10}
Een wiel met straal $30$ cm rolt zonder te slippen. Over welke hoek (in
\omterm{def:g11:trigo:radian}{radialen}, en daarna in graden) draait het wanneer de fiets $1.5$ m
vooruitgaat? Hoever gaat de fiets vooruit bij één volledige omwenteling van
het wiel?
\end{exercise}

\begin{exercise}[$\star\star\star$]\label{exo:g11:trigo:11}
Los op $\intoc{-\pi}{\pi}$ de ongelijkheid $\cos t \leq \frac12$ op. (Schets
de eenheidscirkel, markeer het gebied waar de \omterm{def:g10:coordgeom:system}{abscis} hoogstens $\frac12$ is,
en lees het \omterm{def:g10:numbers:interval}{interval} van waarden van $t$ af.)
\end{exercise}

\section{Opgave: het reuzenrad en het getij}

\begin{problem}[{Weekendopgave --- radialen als afgerolde afstand, en de
eenheidscirkel als hoofdklok van elk periodiek verschijnsel}]
\label{pb:g11:trigo:1}
Een reuzenrad draagt je rond een cirkel; het getij draagt het water van de
haven op en neer langs dezelfde wiskundige cirkel. Elk periodiek
verschijnsel --- wielen, getijden, geluid, seizoenen --- leest zijn
dienstregeling af van de eenheidscirkel uit dit hoofdstuk
(\cref{def:g11:trigo:cossin}), en de \omterm{def:g10:algebra:equation}{vergelijkingen} van
\cref{met:g11:trigo:solve} geven de exacte tijdstippen. Deze opgave rekent
om, rolt op, modelleert en lost op --- en beslecht onderweg waarom
wiskundigen hoeken in \omterm{def:g11:trigo:radian}{radialen} meten.

\medskip
\noindent\textbf{Deel I --- \omterm{def:g11:trigo:radian}{Radialen}, de afgerolde afstand.}
\begin{enumerate}
  \item Zet om naar \omterm{def:g11:trigo:radian}{radialen}: $30^\circ$, $135^\circ$, $300^\circ$; en naar
        graden: $\frac{3\pi}{4}$, $\frac{5\pi}{6}$
        (\cref{met:g11:trigo:convert}).
  \item Waar het bij \omterm{def:g11:trigo:radian}{radialen} om draait: een boog met hoek $t$ (in \omterm{def:g11:trigo:radian}{radialen})
        op een cirkel met straal $r$ heeft lengte precies $r\,t$. Welke boog
        snijdt een hoek van $2$ \omterm{def:g11:trigo:radian}{radialen} af op een cirkel met straal $3$ m?
        (Nergens een $\pi$ --- daar gaat het om.)
  \item Een wiel met straal $30$ cm rolt $1$ km zonder te slippen. Over
        hoeveel \emph{\omterm{def:g11:trigo:radian}{radialen}} is het gedraaid, en hoeveel volledige
        omwentelingen is dat?
  \item Zet uit op de eenheidscirkel en bereken exact:
        $\cos\frac{2\pi}{3}$, $\sin\left(-\frac{\pi}{4}\right)$ en
        $\cos\frac{19\pi}{6}$ (herleid eerst modulo $2\pi$;
        \cref{prop:g11:trigo:associated}).
  \item Een hoek $t \in \intoo{\frac\pi2}{\pi}$ voldoet aan
        $\sin t = \frac35$. Zoek met $\cos^2 t + \sin^2 t = 1$ en het
        kwadrant de waarden van $\cos t$ en $\tan t$.
\end{enumerate}

\medskip
\noindent\textbf{Deel II --- Het reuzenrad.}
Een reuzenrad heeft straal $60$ m, zijn as hangt $65$ m boven de grond, en
het draait één keer rond per $30$ minuten. Je stapt op het laagste punt in op
tijdstip $t = 0$ (in minuten), zodat je hoogte voldoet aan
\[
h(t) = 65 - 60 \cos\!\left(\frac{\pi t}{15}\right).
\]
\begin{enumerate}[resume]
  \item Verantwoord de formule: ga de hoek na waarover na $t$ minuten
        gedraaid is, de hoogte in $t = 0$, en verklaar het minteken.
  \item Bereken je hoogte in $t = 7.5$, $t = 15$ en $t = 20$ minuten (exacte
        waarden).
  \item Op welk tijdstip bevind je je voor het \emph{eerst} op $95$ m? (Los
        $h(t) = 95$ op met \cref{met:g11:trigo:solve}.)
  \item Op welk tijdsinterval van de rit bevind je je op minstens $95$ m
        hoogte --- en hoe lang in totaal? (Vergelijk \cref{exo:g11:trigo:11}.)
  \item Waar verandert de hoogte het \emph{snelst}? Vergelijk de klim tijdens
        de eerste minuut ($h(1) - h(0)$) met de klim tussen minuut $7$ en
        minuut $8$, en formuleer de waarneming (het gereedschap dat haar
        exact maakt --- de \omterm{def:g11:deriv:derivative}{afgeleide} van de sinusoïde --- komt in jaar 12).
  \item De ingenieur aan zet: ontwerp de hoogteformule van een rad met een
        grootste hoogte van $40$ m, een kleinste hoogte van $4$ m en een
        periode van $12$ minuten, waarbij je onderaan instapt in $t = 0$.
\end{enumerate}

\medskip
\noindent\textbf{Deel III --- Het getij en de symmetrieën van de cirkel.}
In een haven bedraagt de waterdiepte in meter, $t$ uur na middernacht,
\[
d(t) = 7 + 3 \sin\!\left(\frac{\pi t}{6}\right).
\]
\begin{enumerate}[resume]
  \item Geef de grootste en de kleinste diepte, de tijdstippen waarop ze
        zich voor het eerst voordoen, en de periode van het getij.
  \item Een vrachtschip heeft $8.5$ m water nodig. Los $d(t) \geq 8.5$ op op
        $\intco{0}{12}$: tijdens welk venster kan het binnenvaren, en hoe
        lang duurt dat venster? Wanneer opent het volgende?
  \item Leid twee identiteiten voor \omterm{prop:g11:trigo:associated}{verwante hoeken} rechtstreeks uit de
        symmetrieën van de cirkel af --- $\sin(\pi - t) = \sin t$ (spiegeling
        om de verticale as) en $\cos(\pi + t) = -\cos t$ (halve draai) --- en
        formuleer de complementaire identiteit
        $\cos\left(\frac\pi2 - t\right) = \sin t$.
  \item Pariteit (met de woordenschat van \cref{pb:g11:func:1}): bepaal of
        $\sin$, $\cos$ en $t \mapsto \sin t + t^3$ even of oneven zijn, met
        bewijzen van één regel vanuit de cirkel.
  \item Los op op $\intco{0}{2\pi}$: $2\sin^2 t - \sin t - 1 = 0$. (Een
        kwadratische \omterm{def:g10:algebra:equation}{vergelijking} in vermomming: ontbind ze zoals in
        \cref{pb:g11:quad:1}, en lees daarna de cirkel af.)
\end{enumerate}

\medskip
\noindent\textbf{Deel IV --- De hoek van de wiskundige.}
\begin{enumerate}[resume]
  \item Hoeveel graden is $1$ \omterm{def:g11:trigo:radian}{radiaal}? En dan een valstrik met de
        rekenmachine: wat is groter, $\sin(1)$ (in \omterm{def:g11:trigo:radian}{radialen}) of
        $\sin(1^\circ)$ --- en ruwweg met welke factor? Welke moraal voor de
        instellingen van je rekenmachine?
  \item Los $\cos t = \sin t$ op op $\intco{0}{2\pi}$, met de cirkel (waar
        zijn \omterm{def:g10:coordgeom:system}{abscis} en \omterm{def:g10:coordgeom:system}{ordinaat} gelijk?).
  \item Voor kleine hoeken vallen de boog en zijn \omterm{def:g11:trigo:cossin}{sinus} bijna samen:
        vergelijk $\sin(0.1)$ met $0.1$ (vijf decimalen), en geef de
        relatieve fout. Die ``kleinehoekbenadering'' stuurt schepen en
        slingers aan; de stelling erachter ($\frac{\sin t}{t} \to 1$) is een
        ster van het limietenhoofdstuk in jaar 12.
  \item Slotstuk --- de eenheidscirkel als hoofdklok, telkens één zin:
        \omterm{def:g11:trigo:radian}{radialen} maken van hoeken eerlijke getallen (boog $=$ straal
        $\times$ hoek); \omterm{def:g11:trigo:cossin}{cosinus} en \omterm{def:g11:trigo:cossin}{sinus} zijn de schaduw en de hoogte van de
        wijzer; de symmetrieën van de cirkel zijn de identiteiten;
        $\cos t = c$ oplossen leest de tijdstippen van de klok af; en elk
        periodiek verschijnsel --- rad, getij, geluid --- is deze klok in een
        kostuum.
\end{enumerate}
\end{problem}
